\documentclass[11pt,a4paper]{article}
\usepackage[utf8]{inputenc}
\usepackage{booktabs}
\usepackage{longtable}
\usepackage{geometry}
\usepackage{hyperref}
\usepackage{amsmath}
\usepackage{amssymb}
\usepackage{graphicx}
\usepackage{float}
\geometry{margin=1in}

% Suppress automatic section numbering (use manual S3, S4, etc.)
\setcounter{secnumdepth}{0}

\begin{document}

\begin{center}
\Large\textbf{Supplementary Materials (Continued)}\\[0.5em]
\normalsize Fact-Based Drug-Drug Interaction Knowledge Graph
\end{center}

\section{S3: Knowledge Graph Construction}

\subsection{S3.1 Data Sources and Integration}

The knowledge graph integrates three authoritative biomedical databases using exact identifier matching only (no fuzzy matching or imputation).

\begin{table}[H]
\centering
\caption{Data source summary.}
\label{tab:sources}
\begin{tabular}{@{}llrl@{}}
\toprule
\textbf{Source} & \textbf{Version} & \textbf{Matched} & \textbf{Contribution} \\
\midrule
DrugBank & 5.1.9 & 4,313 drugs & DDIs, proteins, pathways, categories \\
SIDER & 4.1 & 1,110 drugs & 265,238 drug-side effect pairs \\
CTD & 2024 & 1,999 drugs & 63,278 drug-disease associations \\
\bottomrule
\end{tabular}
\end{table}

\textbf{Matching Methods:}
\begin{itemize}
    \item DrugBank: Primary identifier (DB\#\#\#\#\#)
    \item SIDER: Exact case-insensitive drug name
    \item CTD: CAS number (primary) or exact drug name (fallback)
\end{itemize}

\subsection{S3.2 Knowledge Graph Statistics}

\begin{table}[H]
\centering
\caption{Node statistics.}
\begin{tabular}{@{}lrr@{}}
\toprule
\textbf{Node Type} & \textbf{Count} & \textbf{\%} \\
\midrule
Pathways & 25,958 & 56.9\% \\
Side Effects & 5,548 & 12.2\% \\
Drugs & 4,313 & 9.4\% \\
Categories & 3,619 & 7.9\% \\
Proteins & 3,176 & 7.0\% \\
Diseases & 3,041 & 6.7\% \\
\midrule
\textbf{Total} & \textbf{45,655} & 100\% \\
\bottomrule
\end{tabular}
\end{table}

\begin{table}[H]
\centering
\caption{Edge statistics.}
\begin{tabular}{@{}lrr@{}}
\toprule
\textbf{Edge Type} & \textbf{Count} & \textbf{\%} \\
\midrule
Drug-Drug Interaction & 759,774 & 62.7\% \\
Drug-Side Effect & 265,238 & 21.9\% \\
Drug-Category & 70,618 & 5.8\% \\
Drug-Disease & 63,278 & 5.2\% \\
Drug-Pathway & 31,207 & 2.6\% \\
Drug-Protein & 21,559 & 1.8\% \\
\midrule
\textbf{Total} & \textbf{1,211,674} & 100\% \\
\bottomrule
\end{tabular}
\end{table}

\subsection{S3.3 DDI Severity Distribution}

\begin{table}[H]
\centering
\caption{Recalibrated DDI severity distribution.}
\begin{tabular}{@{}lrrr@{}}
\toprule
\textbf{Severity} & \textbf{Count} & \textbf{\%} & \textbf{DDInter Target} \\
\midrule
Moderate & 608,742 & 80.1\% & 74\% \\
Minor & 70,682 & 9.3\% & 4\% \\
Contraindicated & 44,306 & 5.8\% & 4\% \\
Major & 36,044 & 4.7\% & 18\% \\
\midrule
\textbf{Total} & \textbf{759,774} & 100\% & --- \\
\bottomrule
\end{tabular}
\end{table}

\section{S4: Network Analysis Methods}

\subsection{S4.1 Graph Construction}

The DDI network is an undirected graph $G = (V, E)$ where:
\begin{itemize}
    \item $V$ = 4,314 drug nodes
    \item $E$ = 379,917 undirected edges (unique drug pairs)
\end{itemize}

Note: Raw data contains 759,774 records (bidirectional). The undirected representation counts each pair once.

\subsection{S4.2 Global Network Metrics}

\begin{table}[H]
\centering
\caption{DDI network global metrics.}
\begin{tabular}{@{}llr@{}}
\toprule
\textbf{Metric} & \textbf{Definition} & \textbf{Value} \\
\midrule
Nodes ($|V|$) & Unique drugs & 4,314 \\
Edges ($|E|$) & Unique DDI pairs & 379,917 \\
Density & $\frac{2|E|}{|V|(|V|-1)}$ & 0.041 \\
Average Degree & $\frac{2|E|}{|V|}$ & 176.1 \\
Clustering Coefficient & Global transitivity & 0.70 \\
Average Path Length & Mean shortest path & 2.19 \\
Diameter & Maximum shortest path & 5 \\
Components & Connected subgraphs & 1 \\
\bottomrule
\end{tabular}
\end{table}

\subsection{S4.3 Power-Law Degree Distribution}

The degree distribution follows:
\[
P(k) \propto k^{-\gamma}, \quad \gamma = 0.74, \quad R^2 = 0.71
\]

The low $\gamma$ indicates a ``fat tail'' with many hub drugs having disproportionately high interaction counts ($k_{max} = 2{,}391$).

\subsection{S4.4 Louvain Community Detection}

The Louvain algorithm is a greedy optimization method for detecting community structure in large networks. It operates in two iterative phases:

\textbf{Phase 1 (Local Moving):} Each node is initially assigned to its own community. For each node $i$, the algorithm calculates the modularity gain from moving $i$ to each neighboring community $C$:
\[
\Delta Q = \left[ \frac{\Sigma_{in} + 2k_{i,in}}{2m} - \left( \frac{\Sigma_{tot} + k_i}{2m} \right)^2 \right] - \left[ \frac{\Sigma_{in}}{2m} - \left( \frac{\Sigma_{tot}}{2m} \right)^2 - \left( \frac{k_i}{2m} \right)^2 \right]
\]
where $\Sigma_{in}$ is the sum of weights inside community $C$, $\Sigma_{tot}$ is the total degree of nodes in $C$, $k_i$ is the degree of node $i$, $k_{i,in}$ is the sum of weights from $i$ to nodes in $C$, and $m$ is the total edge weight of the graph.

\textbf{Phase 2 (Aggregation):} Communities are collapsed into super-nodes, creating a new network where edge weights between super-nodes equal the sum of inter-community edges.

These phases repeat until no further modularity improvement is possible, optimizing the global modularity:
\[
Q = \frac{1}{2m} \sum_{ij} \left[ A_{ij} - \frac{k_i k_j}{2m} \right] \delta(c_i, c_j)
\]

\begin{table}[H]
\centering
\caption{Community detection results ($Q = 0.23$).}
\begin{tabular}{@{}lrr@{}}
\toprule
\textbf{Community} & \textbf{Size} & \textbf{\%} \\
\midrule
1 & 1,235 & 28.6\% \\
2 & 858 & 19.9\% \\
3 & 766 & 17.8\% \\
4 & 727 & 16.9\% \\
5 & 708 & 16.4\% \\
6 & 20 & 0.5\% \\
\bottomrule
\end{tabular}
\end{table}

\subsection{S4.5 Centrality Analysis}

\textbf{Degree Centrality:} Measures the number of direct connections a drug has.
\[
C_D(v) = \frac{\text{deg}(v)}{|V| - 1}
\]

\textbf{Betweenness Centrality:} Quantifies how often a drug lies on the shortest path between other drugs in the network.
\[
C_B(v) = \sum_{s \neq v \neq t} \frac{\sigma_{st}(v)}{\sigma_{st}}
\]
where $\sigma_{st}$ is the total number of shortest paths from node $s$ to node $t$, and $\sigma_{st}(v)$ is the number of those paths passing through $v$. Drugs with high betweenness centrality act as ``bridges'' within the interaction network and reflect indirect effects between different drug clusters. Local anesthetics such as procaine and lidocaine, as well as the biologics sotatercept and abciximab, showed high betweenness. These agents interact with multiple drug classes and bridge distinct therapeutic groups.

\textbf{Eigenvector Centrality:} Emphasizes not only the number of connections a drug has but also the importance of its neighboring nodes.
\[
C_E(v) = \frac{1}{\lambda} \sum_{u \in N(v)} C_E(u)
\]
where $\lambda$ is the largest eigenvalue of the adjacency matrix and $N(v)$ is the set of neighbors of $v$. A drug with a high eigenvector score tends to interact with other highly connected drugs. Antiarrhythmics such as quinidine and cardiovascular agents including propranolol, verapamil, and isradipine were at the top of eigenvector centrality rankings, demonstrating their tendency to interact with other high-risk medications.

\textbf{PageRank Centrality:} Measures how likely a drug is to be reached through random traversal of the interaction network.
\[
PR(v) = \frac{1-d}{|V|} + d \sum_{u \in B(v)} \frac{PR(u)}{L(u)}
\]
where $d$ is the damping factor (typically 0.85), $B(v)$ is the set of nodes linking to $v$, and $L(u)$ is the out-degree of $u$. PageRank serves as an indicator of high-interaction targets---drugs most likely to be affected when new drugs are added to a regimen. A drug has high PageRank if many other drugs interact with it, especially if those drugs themselves have many interactions. Local anesthetics (Procaine, Lidocaine, Benzocaine) exhibited the highest PageRank scores, indicating their extensive interaction profiles with multiple drug classes.

\subsection{S4.6 Small-World Properties}

\begin{table}[H]
\centering
\caption{Small-world analysis.}
\begin{tabular}{@{}lrrr@{}}
\toprule
\textbf{Property} & \textbf{DDI Network} & \textbf{Random Graph} & \textbf{Ratio} \\
\midrule
Clustering ($C$) & 0.70 & 0.041 & 17.1$\times$ \\
Path Length ($L$) & 2.19 & 1.76 & 1.25$\times$ \\
\bottomrule
\end{tabular}
\end{table}

The network satisfies small-world criteria: $C \gg C_{random}$ and $L \approx L_{random}$.

\section{S5: Reproducibility}

\subsection{S5.1 Software Dependencies}

\begin{table}[H]
\centering
\caption{Core software requirements.}
\begin{tabular}{@{}lll@{}}
\toprule
\textbf{Library} & \textbf{Version} & \textbf{Purpose} \\
\midrule
Python & 3.10+ & Core language \\
NetworkX & 3.1+ & Graph analysis \\
Pandas & 2.0+ & Data manipulation \\
python-louvain & 0.16+ & Community detection \\
\bottomrule
\end{tabular}
\end{table}

\subsection{S5.2 Execution}

\begin{verbatim}
# Build knowledge graph
python build_fact_based_kg.py

# Network analysis
python ddi_knowledge_graph.py
\end{verbatim}

\subsection{S5.3 Data Access}

\begin{itemize}
    \item DrugBank: \url{https://go.drugbank.com/} (academic license)
    \item SIDER: \url{http://sideeffects.embl.de/} (CC BY-NC-SA)
    \item CTD: \url{http://ctdbase.org/} (open access)
\end{itemize}

\end{document}
